\section{Relocation Model}

\begin{itemize}
  \item VIV-32 v0.1 uses static relocation at link time.
  \item Runtime relocation is not required in v0.1.
  \item Relocations are produced by the assembler and consumed by the linker.
  \item The linker resolves symbols and emits a flat binary image or other executable format.
\end{itemize}

\section{Initial Relocation Requirements}

\begin{itemize}
  \item The initial relocation set should support:
  \begin{itemize}
    \item absolute 32-bit addresses;
    \item PC-relative branch targets;
    \item PC-relative call targets;
    \item split high/low immediate address construction;
    \item section-relative references.
  \end{itemize}
\end{itemize}

\section{Deferred Relocation Features}

\begin{itemize}
  \item Runtime relocation is deferred.
  \item Dynamic linking relocations are deferred.
  \item Global offset table relocations are deferred.
  \item Procedure linkage table relocations are deferred.
  \item Thread-local storage relocations are deferred.
  \item ELF relocation numbers are deferred until a VIV-32 ELF ABI is defined.
\end{itemize}

\section{Position-Independent Code}

\begin{itemize}
  \item Position-independent code is not required in VIV-32 v0.1.
  \item Relocation design should not prevent future PIC support.
  \item Future PIC support may require additional relocations for PC-relative data access and global offset table entries.
\end{itemize}
